\documentclass[a5paper,10pt]{article}

% https://github.com/InDevRus/filippov-solutions

% Нумерация страниц
\pagenumbering{gobble}

% Настройка полей
\usepackage{geometry}
\geometry{left=1cm}
\geometry{right=1cm}
\geometry{top=1cm}
\geometry{bottom=1.5cm}

% Вставка таблицы
\usepackage{array}
\usepackage{tabularx}

% Инструменты выравнивания формул
\usepackage{amsmath}
\usepackage{mathtools}

% Диакритика в математике
\usepackage{amsfonts}

% Вычисления размеров на ходу
\usepackage{calc}

% Удобные записи производных
\usepackage[italic=true]{derivative}

% Настройки сносок
\usepackage{enotez}

% Длинные знаки векторов
\usepackage{anyfontsize}
\usepackage[b]{esvect}

% Вставка картинок
\usepackage{graphicx}

% Вставка красной строки
\usepackage{indentfirst}
\setlength{\parindent}{1em}

% Условия в командах
\usepackage{xifthen}

% Луа-скрипты
\usepackage{luacode}

% Настройка команд отдельно в математическом и текстовом режимах
\usepackage{mathcommand}

% Для повторения LaTeX кода
\usepackage{multido}

% Конфигурация отступов формул
\delimitershortfall-1sp
\usepackage{mleftright}
\mleftright

% Растягивание объектов
\usepackage{relsize}

% Обтекание текстом
\usepackage{wrapfig}
\usepackage{subcaption}
\usepackage{floatflt}

% Поддержка ввода кириллицы
\usepackage[english, russian]{babel}

% Настройка корневой позиции для компилятора
\usepackage{import}
\providecommand*{\ProjectRootPath}{..}

\import{\ProjectRootPath/preambles/macros/}{theorems}

\import{\ProjectRootPath/preambles/fonts}{font_setup}

% Знак градуса
\usepackage{siunitx}

\import{\ProjectRootPath/preambles/macros/}{operators}
\import{\ProjectRootPath/preambles/macros/}{redefinitions}

\import{\ProjectRootPath/preambles/fonts}{letters_setup}
\import{\ProjectRootPath/preambles/fonts/adjustments}{custom_superscripts}

\import{\ProjectRootPath/preambles/custom}{formatting}
\import{\ProjectRootPath/preambles/custom}{frames}
\import{\ProjectRootPath/preambles/custom}{list_setup}

% TODO: перевести команды на современный стиль объявления

\begin{document}

$3\pow{y}{2} y\` + 16x = 2x\pow{y}{3}$.
$y\` = \dfrac {2x\br{\pow{y}{3} - 8}} {3\pow{y}{2}}$.
При этом $y\br{x} = 0$ -- не решение, а $y\br{x} = 2$ является решением.
$\dfrac {3\pow{y}{2} y\`} {\pow{y}{3} - 8} = 2x$.
$\tint \dfrac {3\pow{y}{2} y\`} {\pow{y}{3} - 8} \di x
= \tint \dfrac {3\pow{y}{2} \di y}{\pow{y}{3} - 8}
= \tint \dfrac {1} {\pow{y}{3} - 8} d\br{\pow{y}{3}}
= \ln\vbr{\pow{\y}{3} - 8} + C$.

$\ln\vbr{\pow{\y}{3} - 8} = \pow{x}{2} + C$.
$\pow{y}{3} - 8 = C \pow{e}{\pow{x}{2}}$.
$y\br{x} = \sqrt[3]{C \pow{e}{\pow{x}{2}} + 8}$.
Более удобная форма имеет вид
$y\br{x} = 2 \sqrt[3]{C \pow{e}{\pow{x}{2}} + 1}$,
причём решение $y\br{x} = 2$ будет входить в общее при $C = 0$.

Теперь найдём $y\br{x}$, ограниченные при $x \to +\infty$.

Если $C \ne 0$, то при достаточно больших $x$
$\Big($например, при $x > \sqrt{\max\br{0; -\ln\vbr{C}}}$$\Big)$
имеем $\vbr{C \pow{e}{\pow{x}{2}}} > 1$, а потому
$$\vbr{\y\br{x}}
= 2{\sqrt[3]{\vbr{C \pow{e}{\pow{x}{2}} + 1}}}
\ge 2{\sqrt[3]{\vbr{\vbr{C \pow{e}{\pow{x}{2}}} - 1}}}
= 2{\sqrt[3]{\vbr{C} \pow{e}{\pow{x}{2}} - 1}}
> 2{\sqrt[3]{\vbr{C} \pow{e}{\pow{x}{2}}}}
= 2 \sqrt[3]{\vbr{C}} \exp\br{\dfrac {1} {3} \pow{x}{2}}$$
При $x \to +\infty$ функция 
$g\br{x} = 2 \sqrt[3]{\vbr{C}} \exp\br{\dfrac {1} {3} \pow{x}{2}}$
неограничена, так что неограничена и $y\br{x}$.

Таким образом, единственное ограниченное при 
$x \to +\infty$ решение -- это
$y\br{x} = 2$.
\end{document}
